\subsection{\texorpdfstring{第~\ref{sec:orientability}~节}{可定向性}}\label{subsec:solutions-orientability}

\begin{description}[font=\normalfont,leftmargin=2em,labelwidth=1.5em]
\RGsolution{ex:characterize-orientable-manifold-by-atlas}{characterize-orientable-manifold-by-atlas}
设 $M$ 可定向, 并固定 $M$ 的一个定向. 则对每个 $p\in M$, 都存在 $p$ 附近的参数化 $(U,\varphi)$, 使 $(\mathrm{d}\varphi)_x$ 在每个 $x\in U$ 处保持 $\mathbb{R}^n$ 的标准定向. 取由这些参数化构成的图册 $\mathcal{A}$, 则过渡映射保持定向, 且对每对 $(U_\alpha,\varphi_\alpha),(U_\beta,\varphi_\beta)\in\mathcal{A}$, 在 $W:=\varphi_\alpha(U_\alpha)\cap\varphi_\beta(U_\beta)$ 上有 $\det\mathrm{d}(\varphi_\beta^{-1}\circ\varphi_\alpha)>0$. 反之, 若存在图册 $\mathcal{A}$, 其过渡映射保持定向, 则可如下选定 $M$ 的定向. 对每个 $p\in M$, 取 $p$ 附近的参数化 $(U,\varphi)\in\mathcal{A}$, 并给 $T_pM$ 中与有序基 $\{(\mathrm{d}\varphi)_xe_i\}_{i=1}^n$ 等价的有序基赋予正号, 其中 $x:=\varphi^{-1}(p)$, $\{e_i\}_{i=1}^n$ 是 $\mathbb{R}^n$ 的标准基. 由于所有过渡映射都保持定向, 这个定向的选取与参数化 $\varphi$ 无关.

\RGsolution{ex:orientation-of-bases-along-curve}{orientation-of-bases-along-curve}
设 $t_0\in I$, 并取 $c(t_0)$ 附近的参数化 $(U,\varphi)$, 使得对所有 $x\in U$, $(\mathrm{d}\varphi)_x$ 都保持定向. 在相应的局部坐标 $x:=\varphi^{-1}$ 下, 有
\[
  V_i(t)=\sum_{j=1}^n V_i^j(t)\left(\frac{\partial}{\partial x^j}\right)_{c(t)},
\]
其中函数 $V_i^j:J\to\mathbb{R}$ 定义在 $t_0$ 的邻域 $J\subset I$ 上, 且都是光滑的. 因此由 $d(t)=\det(V_i^j(t))$ 定义的映射 $d:J\to\mathbb{R}$ 也是光滑的. 此外, 由于 $\{V_1(t_0),\ldots,V_n(t_0)\}$ 是 $T_{c(t_0)}M$ 的一组基, 因此或者 $d(t_0)>0$, 或者 $d(t_0)<0$. 从而, 在 $J$ 内 $t_0$ 的某个邻域中, 对 $t$ 也有 $d(t)>0$ 或 $d(t)<0$. 所以使 $\{V_1(t),\ldots,V_n(t)\}$ 为正定向的点 $t\in I$ 构成 $I$ 的开子集, 使它为负定向的点也构成 $I$ 的开子集. 由于 $I$ 连通 (它是一个区间), 可知对所有 $t\in I$, $\{V_1(t),\ldots,V_n(t)\}$ 都是正定向的, 或者都是负定向的.

\RGsolution{ex:mobius-band-is-not-orientable}{mobius-band-is-not-orientable}
设 $c:[0,2\pi]\to M$ 是 Möbius 带 $M$ 上由
\[
  c(\varphi):=g(0,\varphi)=(\cos\varphi,\sin\varphi,0)
\]
定义的曲线, 并考虑沿 $c$ 的光滑向量场, 由
\begin{align*}
  V_1(\varphi)&=\frac{\partial g}{\partial t}(0,\varphi)
  =\left(\cos\frac{\varphi}{2}\cos\varphi,\cos\frac{\varphi}{2}\sin\varphi,\sin\frac{\varphi}{2}\right);\\
  V_2(\varphi)&=\frac{\partial g}{\partial\varphi}(0,\varphi)=(-\sin\varphi,\cos\varphi,0)
\end{align*}
定义. 由于对所有 $\varphi\in[0,2\pi]$,
\[
  \lVert V_1(\varphi)\rVert=\lVert V_2(\varphi)\rVert=1
  \quad\text{且}\quad\langle V_1(\varphi),V_2(\varphi)\rangle=0,
\]
所以对所有 $\varphi\in[0,2\pi]$, 这两个向量场都给出 $T_{c(\varphi)}M$ 的一组基. 又因为 $c(0)=c(2\pi)=(1,0,0)$, 所以
\[
  \{V_1(0),V_2(0)\}=\{(1,0,0);(0,1,0)\}
\]
和
\[
  \{V_1(2\pi),V_2(2\pi)\}=\{(-1,0,0);(0,1,0)\}
\]
是 $T_{(1,0,0)}M$ 的两组基. 然而, 从第一组基到第二组基的换基矩阵
\[
  S=\begin{pmatrix}-1&0\\0&1\end{pmatrix}
\]
具有负的行列式. 因此 $M$ 不可能定向.

\RGsolution{ex:orientable-double-covering}{orientable-double-covering}
\begin{enumerate}[label=(\alph*),ref=\alph*,leftmargin=*]
\item 首先注意, 对 $M$ 的每个参数化 $(U,\varphi)$, 都存在参数化 $(\widetilde U,\widetilde\varphi)$, 它在每个 $p\in\varphi(U)$ 处诱导 $T_pM$ 的相反定向 (只需调换 $\mathbb{R}^n$ 中两个变量的次序).

映射 $\overline{\varphi}:U\to\overline M$ 定义了 $\overline M$ 上的拓扑, 其基为
\[
  \{\overline{\varphi}(U)\mid (U,\varphi)\text{ 是 }M\text{ 的参数化}\}.
\]
这确实是一个拓扑的基, 且此拓扑满足 Hausdorff 公理和第二可数公理, 这是因为
\[
  \{\varphi(U)\mid (U,\varphi)\text{ 是 }M\text{ 的参数化}\}
\]
是 $M$ 上具有同样性质的拓扑的基.

给定 $M$ 的两个参数化 $(U,\varphi)$ 与 $(V,\psi)$, 映射 $\overline{\psi}^{-1}\circ\overline{\varphi}$ 定义在 $\varphi(U)\cap\psi(V)$ 中那些由 $\varphi$ 与 $\psi$ 诱导的定向一致的连通分支在 $\varphi^{-1}$ 下的像上. 因此 $\overline{\psi}^{-1}\circ\overline{\varphi}$ 不仅可微 (在有定义的点处, $\overline{\psi}^{-1}\circ\overline{\varphi}=\psi^{-1}\circ\varphi$), 而且保持定向. 可知
\[
  \{(U,\overline{\varphi})\mid (U,\varphi)\text{ 是 }M\text{ 的参数化}\}
\]
是 $\overline M$ 的一个图册, 其过渡映射保持定向. 因此, $\overline M$ 是一个\textbf{可定向的} $n$ 维流形 [参见 Exercise~\ref{ex:orientability}(\ref{ex:characterize-orientable-manifold-by-atlas})].
\item 这直接来自如下两点: 对上述参数化, $\varphi^{-1}\circ\pi\circ\overline\varphi$ 是恒等映射; 并且如上所述, 对 $M$ 的每个参数化 $(U,\varphi)$, 都存在参数化 $(\widetilde U,\widetilde\varphi)$, 使得对每个 $p\in\varphi(U)$, 它在 $T_pM$ 上诱导相反的定向.
\item 若 $\overline W$ 是 $\overline M$ 的连通分支, 则 $\pi:\overline W\to M$ 仍是覆盖映射. 因此纤维 $\pi^{-1}(p)\subset\overline W$ 或由一点组成, 或由两点组成. 在前一种情形下, $\pi:\overline W\to M$ 是微分同胚, 因而 $M$ 会是可定向的 (因为 $\overline W$ 是可定向的). 事实并非如此, 所以 $\overline W=\overline M$, 从而 $\overline M$ 连通.
\item 恒等式 $\pi\circ\sigma=\pi$ 与 $\sigma\circ\sigma=\sigma$ 立即成立. 由 $\pi\circ\sigma=\pi$ 且 $\pi$ 是局部微分同胚, 可知 $\sigma$ 是微分同胚.
\item 这是因为单连通流形仅容许平凡覆盖映射.
\end{enumerate}
\end{description}
